\ozet{Kubernetes tabanlı sistemlerde güvenlik gözlemlenebilirliği, mikroservisler arası yoğun ağ iletişimi ve kısa ömürlü çalışma yükleri nedeniyle geleneksel yöntemlerle sağlanamamaktadır. Bu tezde, ağ telemetrisi ve çalışma zamanı olaylarının birlikte analizine dayanan çok katmanlı bir anomali tespiti yaklaşımı araştırılmaktadır. Çalışma kapsamında, Kubernetes ortamında eBPF tabanlı Hubble aracı ile ağ akış verileri, Tetragon aracı ile çalışma zamanı olayları toplanmakta; bu veriler ön işleme ve özellik çıkarımı adımlarından geçirilerek katman bazlı gözetimsiz anomali tespit modelleri eğitilmektedir. Ağ ve çalışma zamanı katmanlarından elde edilen anomali skorları, her pod için birleşik bir güvenlik skoru üretmek üzere birleştirilmektedir. Yaklaşımın etkinliği, tek katmanlı modellerin hazır veri kümeleriyle doğrulanması ve çok katmanlı füzyonun kontrollü saldırı senaryolarıyla üretilen veriler üzerinde değerlendirilmesi yoluyla incelenecektir.}

\anahtarkelimeler{Anomali Tespiti, eBPF, Kubernetes, Makine Öğrenmesi, Siber Güvenlik.}